\documentclass[../principal.tex]{subfiles}
\graphicspath{{\subfix{../imagenes/}}}

\begin{document}
\chapter{Corriente}

\section{Definición}

La corriente eléctrica es el flujo de carga eléctrica que pasa por un conductor en un tiempo determinado.

La corriente eléctrica se simboliza con la letra $I$, la carga eléctrica se simboliza con la letra $Q$ y el tiempo se simboliza con la letra $t$.

\begin{equation}
    I = \frac{Q}{t}
\end{equation}



\section{Unidades de medida}

Las unidades de medida de carga eléctrica son el coulomb (\unit{\coulomb}) y de tiempo es el segundo (\unit{\second}).

La unidad de medida de corriente eléctrica entonces es:

\begin{equation}
    \unit{\coulomb\per\second}
\end{equation}

Por conveniencia se usa el ampere (\unit{\ampere}), donde:

\begin{equation}
    \unit{\ampere} = \unit{\coulomb\per\second}
\end{equation}


\end{document}


